\documentclass{beamer}
\usepackage{graphicx,color,amsmath,amssymb,amsfonts,amsbsy,amsthm}
\usepackage{verbatim}
\usetheme{Berlin}

\begin{document}
	
	\begin{frame}
	\title{
		\vfill \hrulefill \vfill
		\begin{center}
			Assured PNT\\
			Challenges for Current Nav Technologies\\
		\end{center}
	}
	\author{Michael L. Carroll}
	\date{\today\\
		\vfill \hrulefill \vspace*{.25in} \vfill
		\begin{tiny}
			\copyright 2024 by Michael L. Carroll
		\end{tiny}
	}
	\maketitle
	\end{frame}

	
	% Slide 1: Introduction to Assured PNT
	\begin{frame}
		\onslide<1->{\frametitle{Introduction to Assured PNT} \framesubtitle{What is Assured PNT?}}
		\begin{itemize}
			\onslide<2->{\item \textbf{Assured Position, Navigation, and Timing (A-PNT):}}
			\begin{itemize}
				\onslide<3->{\item Provides reliable PNT information, even in GPS-degraded or denied environments}
				\onslide<4->{\item Critical for applications requiring high-precision and resilience, such as military, aviation, and critical infrastructure}
			\end{itemize}
			\onslide<5->{\item \textbf{Main Goals of A-PNT}}
			\begin{itemize}
				\onslide<6->{\item Ensuring accuracy and integrity of position, navigation, and timing data}
				\onslide<7->{\item Providing robustness against jamming, spoofing, and interference}
				\onslide<8->{\item Mitigating risks of depending on GPS
				}
				
			\end{itemize}
		\end{itemize}
		\onslide<8->{\center --------------------------------------------------------------------}
	\end{frame}
	
	% Slide 2: Key Technologies in Assured PNT
	\begin{frame}
		\onslide<1->{\frametitle{Key Technologies in Assured PNT} \framesubtitle{Technologies Critical to A-PNT}}
		\begin{itemize}
			\onslide<2->{\item \textbf{Global Navigation Satellite Systems (GNSS):}}
			\begin{itemize}
				\onslide<3->{\item Primary source of PNT data, with anti-jamming and anti-spoofing techniques (e.g., M-Code and SAASM for secure military applications)}
			\end{itemize}
			\onslide<4->{\item \textbf{Inertial Navigation Systems (INS):}}
			\begin{itemize}
				\onslide<5->{\item Provides self-contained PNT information with accelerometers and gyroscopes, allowing operation in GPS-denied environments}
			\end{itemize}
			\onslide<6->{\item \textbf{Alternative PNT Sources:}}
			\begin{itemize}
				\onslide<7->{\item Terrestrial-based signals (e.g., radio, cellular networks) and signals from Low Earth Orbit (LEO) satellites}
			\end{itemize}
			\onslide<8->{\item \textbf{Precision Timing Sources:}}
			\begin{itemize}
				\onslide<9->{\item Atomic clocks and GPS-disciplined oscillators to ensure reliable timing}
			\end{itemize}
		\end{itemize}
		\onslide<10->{\center --------------------------------------------------------------------}
	\end{frame}
\begin{frame}
	\onslide<1->{\frametitle{Security of M-Code Signals} \framesubtitle{Spoofing and Adversarial Threats}}
	\begin{itemize}
		\onslide<2->{\item \textbf{Purpose of M-Code:}}
		\begin{itemize}
			\onslide<3->{\item Developed by the U.S. Department of Defense (DoD) for enhanced security in military GPS}
			\onslide<4->{\item Provides stronger encryption, anti-jamming, and anti-spoofing features compared to earlier GPS signals}
		\end{itemize}
		\onslide<5->{\item \textbf{Current Status:}}
		\begin{itemize}
			\onslide<6->{\item No publicly available evidence that adversaries have successfully spoofed M-Code signals}
			\onslide<7->{\item M-Code’s design makes spoofing significantly more challenging than with civilian signals}
		\end{itemize}
		\onslide<8->{\item \textbf{DoD Countermeasures:}}
		\begin{itemize}
			\onslide<9->{\item Ongoing assessment and improvements in GPS security to counter emerging threats}
			\onslide<10->{\item Continued vigilance as part of broader efforts to secure military navigation and timing}
		\end{itemize}
	\end{itemize}
	\onslide<11->{\center --------------------------------------------------------------------}
\end{frame}
\begin{frame}
	\onslide<1->{\frametitle{M-Code Resistance to Jamming} \framesubtitle{Key Features Enhancing M-Code Resilience}}
	\begin{itemize}
		\onslide<2->{\item \textbf{Higher Power Signal:}}
		\begin{itemize}
			\onslide<3->{\item M-Code is broadcast at a higher power level than civilian GPS signals, making it harder to overwhelm with jamming signals}
		\end{itemize}
		\onslide<4->{\item \textbf{Advanced Modulation Techniques:}}
		\begin{itemize}
			\onslide<5->{\item Uses Binary Offset Carrier (BOC) modulation, which improves signal separation and reduces interference}
		\end{itemize}
		\onslide<6->{\item \textbf{Narrower Beam Forming:}}
		\begin{itemize}
			\onslide<7->{\item Directed beam shaping on modern GPS III satellites focuses signals for military users, increasing jamming resistance}
		\end{itemize}
		\onslide<8->{\item \textbf{Enhanced Encryption and Anti-Spoofing:}}
		\begin{itemize}
			\onslide<9->{\item Strong encryption and authentication make M-Code signals harder for adversaries to disrupt or imitate}
		\end{itemize}
	\end{itemize}
	\onslide<10->{\center --------------------------------------------------------------------}
\end{frame}
	
	% Slide 3: Achieving Assured PNT
	\begin{frame}
		\onslide<1->{\frametitle{Achieving Assured PNT} \framesubtitle{Strategies for Resilience and Accuracy}}
		\begin{itemize}
			\onslide<2->{\item \textbf{Sensor Fusion Algorithms:}}
			\begin{itemize}
				\onslide<3->{\item Integrate data from multiple sources (GNSS, INS, timing devices) to create a reliable composite solution}
				\onslide<4->{\item Kalman filters and machine learning models are commonly used}
			\end{itemize}
			\onslide<5->{\item \textbf{Robust Signal Authentication:}}
			\begin{itemize}
				\onslide<6->{\item Techniques to verify signal integrity and authenticity, essential for detecting spoofing and interference}
			\end{itemize}
			\onslide<7->{\item \textbf{System Redundancy:}}
			\begin{itemize}
				\onslide<8->{\item Use of alternative PNT sources to maintain accuracy in contested or challenging environments}
			\end{itemize}
			\onslide<9->{\item \textbf{Intermittent GPS Synchronization:}}
			\begin{itemize}
				\onslide<10->{\item Periodic GPS updates improve INS drift performance and timing accuracy}
			\end{itemize}
		\end{itemize}
		\onslide<11->{\center --------------------------------------------------------------------}
	\end{frame}
	
	



	% Slide 7: Assured PNT System Overview
\begin{frame}
	\onslide<1->{\frametitle{Assured Position, Navigation, and Timing (A-PNT)} \framesubtitle{System Components}}
	\begin{itemize}
		\onslide<2->{\item \textbf{Primary Navigation Sources}}
		\begin{itemize}
			\onslide<3->{\item GNSS (e.g., GPS), supplemented with anti-jamming technology}
		\end{itemize}
		\onslide<4->{\item \textbf{Alternative PNT Sources}}
		\begin{itemize}
			\onslide<5->{\item Terrestrial radio, cellular, and other satellite constellations}
		\end{itemize}
		\onslide<6->{\item \textbf{IMUs and Precise Clocks}}
		\begin{itemize}
			\onslide<7->{\item INS data for position and orientation in GPS-denied conditions}
		\end{itemize}
	\end{itemize}
	\onslide<8->{\center --------------------------------------------------------------------}
\end{frame}

% Slide 8: State Model for INS/IMU, GPS, and Timing Kalman Filter
\begin{frame}
	\onslide<1->{\frametitle{Kalman Filter State Model for INS/IMU, GPS, and Timing} \framesubtitle{State Vector}}
	\begin{itemize}
		\onslide<2->{\item Minimal State Vector 
			\footnotesize{\(\mathbf{x} = [P_x, P_y, P_z, V_x, V_y, V_z, \newline
				\text{Orientation}, \text{Accel Bias}, \text{Gyro Bias}, \text{Clock Bias}, \text{Clock Drift}]\)}}
		\onslide<3->{\item \textbf{Position and Velocity}}
		\onslide<4->{\item \textbf{Attitude and Orientation}}
		\onslide<5->{\item \textbf{Accelerometer and Gyroscope Biases}}
		\onslide<6->{\item \textbf{Clock Bias and Drift}}
	\end{itemize}
	\onslide<7->{\center --------------------------------------------------------------------}
\end{frame}

% Slide 1: Introduction to the International GNSS Service
\begin{frame}
	\onslide<1->{\frametitle{Additional Navaids:  The International GNSS Service (IGS)} \framesubtitle{Overview}}
	\begin{itemize}
		\onslide<2->{\item \textbf{What is IGS?}}
		\begin{itemize}
			\onslide<3->{\item IGS is a global collaborative network led by NASA's Jet Propulsion Laboratory (JPL)}
			\onslide<4->{\item Provides highly accurate GNSS data and products, supporting research, navigation, and Earth observation}
			\onslide<5->{\item Supported by 500+ Tracking Stations and Dozens of Analysis Centers}
		\end{itemize}
		\onslide<6->{\item \textbf{GNSS Systems Supported:}}
		\begin{itemize}
			\onslide<7->{\item IGS includes data from multiple GNSS systems: GPS, GLONASS, Galileo, BeiDou, and QZSS}
			\onslide<8->{\item Enables multi-constellation support, increasing resilience when GPS signals are unavailable}
		\end{itemize}
	\end{itemize}
	\onslide<9->{\center --------------------------------------------------------------------}
\end{frame}

% Slide 2: Key IGS Products and Applications
\begin{frame}
	\onslide<1->{\frametitle{Key IGS Products} \framesubtitle{Applications in Navigation and PNT Systems}}
	\begin{itemize}
		\onslide<2->{\item \textbf{Precise GNSS Satellite Orbits:}}
		\begin{itemize}
			\onslide<3->{\item Provides satellite position data with centimeter-level accuracy}
		\end{itemize}
		\onslide<4->{\item \textbf{Satellite Clock Corrections:}}
		\begin{itemize}
			\onslide<5->{\item Offers high-accuracy clock corrections, improving timing accuracy across systems}
		\end{itemize}
		\onslide<6->{\item \textbf{Earth Rotation Parameters (ERP):}}
		\begin{itemize}
			\onslide<7->{\item Supports high-precision geospatial applications by modeling Earth’s rotational variations}
		\end{itemize}
		\onslide<8->{\item \textbf{Ionospheric and Tropospheric Data:}}
		\begin{itemize}
			\onslide<9->{\item Models atmospheric effects on GNSS signals, enabling better correction and positioning accuracy}
		\end{itemize}
	\end{itemize}
	\onslide<10->{\center --------------------------------------------------------------------}
\end{frame}
\begin{frame}
	\onslide<1->{\frametitle{IGS Real-Time Satellite and Clock Corrections} \framesubtitle{IGS Real-Time Service (RTS)}}
	\begin{itemize}
		\onslide<2->{\item \textbf{Real-Time Service (RTS):}}
		\begin{itemize}
			\onslide<3->{\item Provides real-time GPS satellite orbit and clock correction parameters}
			\onslide<4->{\item Distributed via RTCM State Space Representation (SSR) correction streams}
		\end{itemize}
		\onslide<5->{\item \textbf{Supports Multi-Constellation GNSS:}}
		\begin{itemize}
			\onslide<6->{\item Currently includes GPS and GLONASS, with plans for additional GNSS constellations}
		\end{itemize}
		\onslide<7->{\item \textbf{Applications:}}
		\begin{itemize}
			\onslide<8->{\item Enables real-time Precise Point Positioning (PPP) for high-accuracy positioning}
			\onslide<9->{\item Supports scientific, educational, and commercial applications needing real-time corrections}
		\end{itemize}
		\onslide<10->{\item \textbf{Access via NTRIP:}}
		\begin{itemize}
			\onslide<11->{\item Data is accessible globally through IGS NTRIP broadcasters, facilitating real-time use}
		\end{itemize}
	\end{itemize}
	\onslide<12->{\center --------------------------------------------------------------------}
\end{frame}

% Slide 3: Using IGS Products to Bound Gyro Errors in GPS-Denied Environments
\begin{frame}
	\onslide<1->{\frametitle{Bounding Gyro Errors with IGS Products} \framesubtitle{Role of IGS in GPS-Denied Scenarios}}
	\begin{itemize}
		\onslide<2->{\item \textbf{Multi-Constellation GNSS Data:}}
		\begin{itemize}
			\onslide<3->{\item IGS provides real-time GNSS data from systems like GLONASS and Galileo, which can supplement GPS-denied times}
		\end{itemize}
		\onslide<4->{\item \textbf{Leveraging Precise Clock Corrections:}}
		\begin{itemize}
			\onslide<5->{\item Accurate timing corrections from IGS improve clock drift estimates, supporting INS and gyro-based dead reckoning}
		\end{itemize}
		\onslide<6->{\item \textbf{Ionospheric Models for Error Mitigation:}}
		\begin{itemize}
			\onslide<7->{\item IGS’s atmospheric data can be used to refine position estimates, reducing gyro error propagation}
		\end{itemize}
		\onslide<8->{\item \textbf{Earth Rotation Parameters (ERP):}}
		\begin{itemize}
			\onslide<9->{\item ERPs help maintain geodetic accuracy, which stabilizes navigation in the absence of direct GNSS signals}
		\end{itemize}
	\end{itemize}
	\onslide<10->{\center --------------------------------------------------------------------}
\end{frame}

% Slide 4: Benefits of IGS in Resilient PNT Solutions
\begin{frame}
	\onslide<1->{\frametitle{Benefits of IGS in Resilient PNT} \framesubtitle{Enhancing PNT Robustness with IGS Data}}
	\begin{itemize}
		\onslide<2->{\item \textbf{Increased Redundancy through Multi-GNSS Data:}}
		\begin{itemize}
			\onslide<3->{\item Enables fallback to alternative GNSS systems (e.g., GLONASS, Galileo), maintaining PNT resilience}
		\end{itemize}
		\onslide<4->{\item \textbf{Reduced Dependence on Single GNSS System:}}
		\begin{itemize}
			\onslide<5->{\item IGS provides products compatible with multiple constellations, lowering the impact of GPS outages}
		\end{itemize}
		\onslide<6->{\item \textbf{Supports Continuous Navigation:}}
		\begin{itemize}
			\onslide<7->{\item By integrating IGS products, INS and gyro systems maintain bounded errors and reliable navigation}
		\end{itemize}
	\end{itemize}
	\onslide<8->{\center --------------------------------------------------------------------}
\end{frame}

	
	% Slide 1: Navigation Aids in GPS-Denied Environments
	\begin{frame}
		\onslide<1->{\frametitle{Navigation Aids for Aircraft in GPS-Denied Environments} \framesubtitle{Assisting INS/ Kalman Filter Systems}}
		\begin{itemize}
			\onslide<2->{\item \textbf{Inertial Navigation System (INS):}}
			\begin{itemize}
				\onslide<3->{\item Primary source when GPS is denied, but accuracy degrades over time}
				\onslide<4->{\item External measurement sources help to bound INS errors}
			\end{itemize}
			\onslide<5->{\item \textbf{Ground-Based Navigation Aids:}}
			\begin{itemize}
				\onslide<6->{\item \textbf{VOR (Very High Frequency Omnidirectional Range):} Provides azimuth or bearing information}
				\onslide<7->{\item \textbf{DME (Distance Measuring Equipment):} Measures slant range for triangulation when combined with VOR}
				\onslide<8->{\item \textbf{TACAN (Tactical Air Navigation):} Combines azimuth and distance, offering precise positioning for military aircraft}
			\end{itemize}
		\end{itemize}
		\onslide<9->{\center --------------------------------------------------------------------}
	\end{frame}
	
	% Slide 2: Additional Ground-Based and Radar-Based Systems
	\begin{frame}
		\onslide<1->{\frametitle{Additional Ground-Based and Radar-Based Aids} \framesubtitle{Supporting Position and Bearing Data}}
		\begin{itemize}
			\onslide<2->{\item \textbf{Automatic Direction Finder (ADF):}}
			\begin{itemize}
				\onslide<3->{\item Uses Non-Directional Beacons (NDBs) for bearing; a backup to VOR/DME}
			\end{itemize}
			\onslide<4->{\item \textbf{Radar-Based Systems:}}
			\begin{itemize}
				\onslide<5->{\item \textbf{Secondary Surveillance Radar (SSR):} Provides range and bearing within controlled airspace}
				\onslide<6->{\item \textbf{Multilateration (MLAT):} Uses transponder signal timing for precise position triangulation}
				\onslide<7->{\item \textbf{Weather Radar:} Assists in dead reckoning using terrain and storm positioning}
			\end{itemize}
		\end{itemize}
		\onslide<8->{\center --------------------------------------------------------------------}
	\end{frame}
	
	% Slide 3: Alternative Satellite and Sensor-Based Aids
	\begin{frame}
		\onslide<1->{\frametitle{Alternative Satellite and Sensor-Based Navigation Aids} \framesubtitle{Augmenting INS Accuracy}}
		\begin{itemize}
			\onslide<2->{\item \textbf{Alternative Satellite Constellations:}}
			\begin{itemize}
				\onslide<3->{\item Multi-GNSS systems (e.g., GLONASS, Galileo) act as backups when GPS is unavailable}
			\end{itemize}
			\onslide<4->{\item \textbf{Onboard Vision and Sensor Systems:}}
			\begin{itemize}
				\onslide<5->{\item \textbf{Terrain-Referenced Navigation (TRN):} Matches terrain features with maps}
				\onslide<6->{\item \textbf{Camera-Based Navigation:} Uses visual cues for position estimation (e.g., visual odometry)}
				\onslide<7->{\item \textbf{Air Data Systems:} Provides wind and airspeed data for dead reckoning}
			\end{itemize}
		\end{itemize}
		\onslide<8->{\center --------------------------------------------------------------------}
	\end{frame}

	
	% Slide 1: Introduction to Terrain-Referenced Navigation (TRN)
	\begin{frame}
		\onslide<1->{\frametitle{Terrain-Referenced Navigation (TRN)} \framesubtitle{Overview}}
		\begin{itemize}
			\onslide<2->{\item \textbf{What is TRN?}}
			\begin{itemize}
				\onslide<3->{\item A navigation technique that uses terrain features to estimate position by matching observed data with known maps}
			\end{itemize}
			\onslide<4->{\item \textbf{Primary Application:}}
			\begin{itemize}
				\onslide<5->{\item Used in GPS-denied situations to provide navigation data in various terrains (e.g., mountains, valleys, coastlines)}
			\end{itemize}
			\onslide<6->{\item \textbf{Common Use Cases:}}
			\begin{itemize}
				\onslide<7->{\item Military aircraft and submarines for navigation in hostile or GPS-denied areas}
				\onslide<8->{\item Autonomous vehicles for precise navigation over land}
			\end{itemize}
		\end{itemize}
		\onslide<9->{\center --------------------------------------------------------------------}
	\end{frame}
	
	% Slide 2: How TRN Works
	\begin{frame}
		\onslide<1->{\frametitle{How TRN Works} \framesubtitle{Matching Terrain Features to Maps}}
		\begin{itemize}
			\onslide<2->{\item \textbf{Step 1: Data Collection}}
			\begin{itemize}
				\onslide<3->{\item Onboard sensors (e.g., radar, LiDAR, or sonar) collect data on terrain features below the vehicle}
			\end{itemize}
			\onslide<4->{\item \textbf{Step 2: Terrain Matching}}
			\begin{itemize}
				\onslide<5->{\item The collected data is matched to a stored digital elevation model (DEM) of the terrain}
				\onslide<6->{\item Algorithms compare observed terrain patterns with the DEM to estimate location}
			\end{itemize}
			\onslide<7->{\item \textbf{Step 3: Position Update}}
			\begin{itemize}
				\onslide<8->{\item Position estimate is continuously updated as the vehicle moves, using successive matches to terrain}
			\end{itemize}
		\end{itemize}
		\onslide<9->{\center --------------------------------------------------------------------}
	\end{frame}
	
	% Slide 3: TRN Sensors and Data Sources
	\begin{frame}
		\onslide<1->{\frametitle{TRN Sensors and Data Sources} \framesubtitle{Onboard Sensors for Data Collection}}
		\begin{itemize}
			\onslide<2->{\item \textbf{Radar Altimeters:}}
			\begin{itemize}
				\onslide<3->{\item Measures distance to the ground, commonly used on aircraft for altitude and terrain mapping}
			\end{itemize}
			\onslide<4->{\item \textbf{LiDAR (Light Detection and Ranging):}}
			\begin{itemize}
				\onslide<5->{\item Uses laser pulses to create detailed terrain profiles; effective for both land and low-altitude applications}
			\end{itemize}
			\onslide<6->{\item \textbf{Sonar (Sound Navigation and Ranging):}}
			\begin{itemize}
				\onslide<7->{\item Used primarily by submarines to map the seafloor or underwater terrain}
			\end{itemize}
		\end{itemize}
		\onslide<8->{\center --------------------------------------------------------------------}
	\end{frame}
	
	% Slide 4: TRN Accuracy and Error Bounds
	\begin{frame}
		\onslide<1->{\frametitle{TRN Accuracy and Error Bounds} \framesubtitle{Factors Influencing Accuracy}}
		\begin{itemize}
			\onslide<2->{\item \textbf{Resolution of Digital Elevation Model (DEM):}}
			\begin{itemize}
				\onslide<3->{\item Higher resolution DEMs enable more accurate terrain matching and smaller error bounds}
				\onslide<4->{\item Typical accuracy with high-resolution DEMs can be within 10-20 meters in well-mapped areas}
			\end{itemize}
			\onslide<5->{\item \textbf{Terrain Variability:}}
			\begin{itemize}
				\onslide<6->{\item Accuracy is higher in areas with distinct terrain features (e.g., mountains, valleys)}
				\onslide<7->{\item Flat or homogeneous terrain may lead to larger error bounds due to lack of unique features}
			\end{itemize}
			\onslide<8->{\item \textbf{Sensor Accuracy and Noise:}}
			\begin{itemize}
				\onslide<9->{\item Higher-quality sensors (e.g., LiDAR vs. radar) generally provide better terrain mapping accuracy}
				\onslide<10->{\item Sensor noise can lead to mismatches with the DEM, increasing position error}
			\end{itemize}
		\end{itemize}
		\onslide<11->{\center --------------------------------------------------------------------}
	\end{frame}
	
	% Slide 5: Limitations and Challenges of TRN
	\begin{frame}
		\onslide<1->{\frametitle{Limitations and Challenges of TRN} \framesubtitle{Considerations for Effective Use}}
		\begin{itemize}
			\onslide<2->{\item \textbf{Dependence on DEM Quality and Availability:}}
			\begin{itemize}
				\onslide<3->{\item High-resolution DEMs may not be available for all regions, reducing TRN effectiveness in those areas}
			\end{itemize}
			\onslide<4->{\item \textbf{Environmental Constraints:}}
			\begin{itemize}
				\onslide<5->{\item Poor weather, low visibility, or extreme terrain (e.g., deserts) can impact sensor readings}
			\end{itemize}
			\onslide<6->{\item \textbf{Resource Requirements:}}
			\begin{itemize}
				\onslide<7->{\item High-performance computing resources needed for real-time terrain matching and updating}
			\end{itemize}
		\end{itemize}
		\onslide<8->{\center --------------------------------------------------------------------}
	\end{frame}
	
	
	% Slide 4: Communication-Based Positioning and RTK
	\begin{frame}
		\onslide<1->{\frametitle{Communication-Based Positioning and RTK} \framesubtitle{Additional Measurement Sources}}
		\begin{itemize}
			\onslide<2->{\item \textbf{Real-Time Kinematic (RTK) Corrections:}}
			\begin{itemize}
				\onslide<3->{\item Ground-based RTK stations provide high-precision correction data, useful with partial GNSS}
			\end{itemize}
			\onslide<4->{\item \textbf{Communication-Based Positioning:}}
			\begin{itemize}
				\onslide<5->{\item \textbf{ADS-B In/Out:} Provides positional data from nearby aircraft, allowing relative positioning}
				\onslide<6->{\item \textbf{Cellular Networks:} Can support triangulation at low altitudes (e.g., over populated areas)}
			\end{itemize}
		\end{itemize}
		\onslide<7->{\center --------------------------------------------------------------------}
	\end{frame}
	
	% Slide 5: Integration with INS/GPS Kalman Filter
	\begin{frame}
		\onslide<1->{\frametitle{Integration with INS/GPS Kalman Filter} \framesubtitle{Enhancing Accuracy in GPS-Denied Conditions}}
		\begin{itemize}
			\onslide<2->{\item \textbf{Kalman Filter Integration:}}
			\begin{itemize}
				\onslide<3->{\item Each alternative navigation aid can provide updates to the Kalman filter}
				\onslide<4->{\item Measurement sources correct INS drift, helping maintain accurate estimates of position, velocity, and attitude}
			\end{itemize}
			\onslide<5->{\item \textbf{Resilience in GPS-Denied Environments:}}
			\begin{itemize}
				\onslide<6->{\item Kalman filter dynamically uses available sources to ensure reliable navigation until GPS is restored}
			\end{itemize}
		\end{itemize}
		\onslide<7->{\center --------------------------------------------------------------------}
	\end{frame}



	% Slide 1: Budgetary Challenges in Achieving Assured PNT
\begin{frame}
	\onslide<1->{\frametitle{Backup:  Budgetary Challenges in Achieving Assured PNT} \framesubtitle{Impact of Program Delays}}
	\begin{itemize}
		\onslide<2->{\item \textbf{Delayed GPS OCX and MGUE Programs:}}
		\begin{itemize}
			\onslide<3->{\item \textbf{GPS OCX (Next Generation Operational Control System):} \newline 
				Critical for improved anti-jamming and anti-spoofing capabilities, but delayed deployment affects A-PNT timelines}
			\onslide<4->{\item \textbf{GPS MGUE (Military GPS User Equipment):} \newline 
				Essential for integrating M-Code, which is designed to enhance signal security; delays impact military readiness for secure PNT}
		\end{itemize}
		\onslide<5->{\item \textbf{Increased Costs Due to Extended Timelines:}}
		\begin{itemize}
			\onslide<6->{\item Project delays require additional funding to maintain and upgrade legacy systems as a stopgap solution}
		\end{itemize}
		\onslide<7->{\item \textbf{Dependency on Supplemental PNT Solutions:}}
		\begin{itemize}
			\onslide<8->{\item Delays in OCX and MGUE push agencies to invest in alternative PNT solutions (e.g., terrestrial or LEO-based systems), increasing overall costs}
		\end{itemize}
	\end{itemize}
	\onslide<9->{\center --------------------------------------------------------------------}
\end{frame}

% Slide 2: Additional Budgetary Considerations
\begin{frame}
	\onslide<1->{\frametitle{Additional Budgetary Considerations for A-PNT} \framesubtitle{Adapting to Delays and Rising Costs}}
	\begin{itemize}
		\onslide<2->{\item \textbf{Balancing Investment in Legacy Systems and Modernization:}}
		\begin{itemize}
			\onslide<3->{\item Funding needed to sustain legacy GPS systems until OCX and MGUE are operational, adding to budget strain}
			\onslide<4->{\item Resources diverted to address immediate operational gaps rather than long-term A-PNT advancements}
		\end{itemize}
		\onslide<5->{\item \textbf{Investment in Alternative A-PNT Technologies:}}
		\begin{itemize}
			\onslide<6->{\item Increased allocation for non-GPS-dependent solutions (e.g., IMUs, precision timing, and LEO constellations) adds to the budget but enhances resilience}
		\end{itemize}
		\onslide<7->{\item \textbf{Cost of Integrating New Technologies:}}
		\begin{itemize}
			\onslide<8->{\item Integrating supplemental technologies with existing infrastructure can be complex and costly}
			\onslide<9->{\item Ensuring interoperability between legacy and new systems is resource-intensive}
		\end{itemize}
	\end{itemize}
	\onslide<10->{\center --------------------------------------------------------------------}
\end{frame}

	
	% Slide 1: Gyroscope Technology Overview
	\begin{frame}
		\onslide<1->{\frametitle{State of the Art in Gyroscope Technology} \framesubtitle{Overview}}
		\begin{itemize}
			\onslide<2->{\item \textbf{Optical Gyroscopes}}
			\begin{itemize}
				\onslide<3->{\item Fiber-Optic Gyroscopes (FOGs) \newline
					Achieve accuracy of 0.1 to 0.0001 degrees per hour}
				\onslide<4->{\item Optical Resonator Gyroscopes (ORGs) \newline
					Random drift values as low as 0.06 degrees per hour}
			\end{itemize}
			\onslide<5->{\item \textbf{Micro-Electro-Mechanical Systems (MEMS) Gyroscopes}}
			\begin{itemize}
				\onslide<6->{\item Miniaturization advances enable high-precision navigation-grade applications}
			\end{itemize}
			\onslide<7->{\item \textbf{Quantum Gyroscopes}}
			\begin{itemize}
				\onslide<8->{\item Using double-mode surface-acoustic-wave cavities to surpass standard quantum limits}
			\end{itemize}
			\onslide<9->{\item \textbf{Applications} \newline Expanding to medical, sensing, imaging, and critical navigation}
		\end{itemize}
		\onslide<10->{\center --------------------------------------------------------------------}
	\end{frame}
	
	% Slide 2: Navigation-Grade Performance Criteria
	\begin{frame}
		\onslide<1->{\frametitle{Navigation-Grade Gyroscope Performance} \framesubtitle{Criteria}}
		\begin{itemize}
			\onslide<2->{\item Bias Instability (Drift Rate)}
			\begin{itemize}
				\onslide<3->{\item Drift rates around 0.01 to 1 degree per hour}
			\end{itemize}
			\onslide<4->{\item Scale Factor Accuracy}
			\begin{itemize}
				\onslide<5->{\item Accuracy within 50 ppm}
			\end{itemize}
			\onslide<6->{\item Noise Performance (Angle Random Walk)}
			\begin{itemize}
				\onslide<7->{\item Below 0.01 degrees per \(\sqrt{\text{hour}}\)}
			\end{itemize}
			\onslide<8->{\item Thermal Stability and Bandwidth}
			\begin{itemize}
				\onslide<9->{\item Consistent performance under varying temperatures and dynamic changes}
			\end{itemize}
		\end{itemize}
		\onslide<10->{\center --------------------------------------------------------------------}
	\end{frame}
	
	% Slide 3: Navigation-Grade MEMS Gyroscopes
	\begin{frame}
		\onslide<1->{\frametitle{MEMS Gyroscopes with Navigation-Grade Performance} \framesubtitle{Examples}}
		\begin{itemize}
			\onslide<2->{\item \textbf{CEA-Leti's Nano-Resistive Sensing Gyroscope}}
			\begin{itemize}
				\onslide<3->{\item Bias instability of 0.02°/h and angular random walk (ARW) of 0.004°/\(\sqrt{\text{h}}\)}
			\end{itemize}
			\onslide<4->{\item \textbf{Safran's Si-MEMS Gyroscope}}
			\begin{itemize}
				\onslide<5->{\item High-precision with optimization for CSWAP requirements}
			\end{itemize}
			\onslide<6->{\item \textbf{EMCORE's Quartz MEMS IMU}}
			\begin{itemize}
				\onslide<7->{\item Angle random walk < 0.006°/\(\sqrt{\text{h}}\), bias stability of 0.2°/h over temperature}
			\end{itemize}
		\end{itemize}
		\onslide<8->{\center --------------------------------------------------------------------}
	\end{frame}
	
	% Slide 4: Role of Precise Clocks in MEMS Gyro-CSWAP Systems
	\begin{frame}
		\onslide<1->{\frametitle{Precise Clocks in MEMS Gyro-CSWAP Systems} \framesubtitle{Importance without GPS}}
		\begin{itemize}
			\onslide<2->{\item \textbf{Reducing Drift Accumulation}}
			\begin{itemize}
				\onslide<3->{\item Synchronizes timing, reducing drift over time}
			\end{itemize}
			\onslide<4->{\item \textbf{Improving Dead Reckoning Accuracy}}
			\begin{itemize}
				\onslide<5->{\item Enables consistent time-stamping of measurements}
			\end{itemize}
			\onslide<6->{\item \textbf{Supporting Sensor Fusion Algorithms}}
			\begin{itemize}
				\onslide<7->{\item Accurate timing of multi-sensor data, improving algorithm performance}
			\end{itemize}
			\onslide<8->{\item \textbf{Supporting Assisted Inertial Navigation Systems (AINS)}}
			\begin{itemize}
				\onslide<9->{\item Assists in periodic GPS syncing for extended stability}
			\end{itemize}
		\end{itemize}
		\onslide<10->{\center --------------------------------------------------------------------}
	\end{frame}
	
	% Slide 5: EMCORE Overview
	\begin{frame}
		\onslide<1->{\frametitle{About EMCORE Corporation} \framesubtitle{Overview and History}}
		\begin{itemize}
			\onslide<2->{\item Founded in 1984}
			\onslide<3->{\item Specializes in fiber optic, MEMS, and ring laser gyro technology}
			\onslide<4->{\item Applications in aerospace, defense, and industrial markets}
		\end{itemize}
		\onslide<5->{\center --------------------------------------------------------------------}
	\end{frame}
	
	% Slide 6: FEI-Zyfer and Assured PNT Applications
	\begin{frame}
		\onslide<1->{\frametitle{FEI-Zyfer and Assured PNT} \framesubtitle{Timing Solutions for PNT}}
		\begin{itemize}
			\onslide<2->{\item \textbf{GPS-Disciplined Time and Frequency Systems}}
			\begin{itemize}
				\onslide<3->{\item M-Code and SAASM GPS receivers for secure applications}
			\end{itemize}
			\onslide<4->{\item \textbf{High-Precision Oscillators}}
			\begin{itemize}
				\onslide<5->{\item Uses Rubidium Atomic Clocks for precise timing}
			\end{itemize}
			\onslide<6->{\item \textbf{Certified for Defense Applications}}
			\begin{itemize}
				\onslide<7->{\item Approved Host Application Equipment integrator}
			\end{itemize}
		\end{itemize}
		\onslide<8->{\center --------------------------------------------------------------------}
	\end{frame}
	
	
\end{document}
